\chapter{はじめに}

\emph{ランタンの芯（wick）は、炎を支える部分です。} wick は lantern エンジン専用のスクリプト言語です。Lua の小ささと使い心地を受け継ぎながら、つまずきやすい部分を lint ルールで補うのではなく、言語そのものの設計で取り除いています。本章では、ゲームエンジンが独自の言語を持つに至った理由、その言語が解決しようとする問題、そして本書の構成を説明します。

\section{wick が生まれた理由}

wick を作ったのは、世の中にもう一つ言語が必要だったからではありません。Lua で実際にゲームを出荷し、困ったことを記録していった結果、どの項目も同じ構図だと気づいたからです。言語が許してしまう誤りがあり、それが見つかるのはプレイヤーが遭遇したときだけでした。

Lua はすばらしい言語です。小さく、組み込みやすく、高速で、C との連携もとても簡単です。この十年、「ゲームのスクリプトには何を使うべきか」という問いへの定番の答えでした。その長所に異を唱えるつもりはありません。ただ、ゲームエンジンのスクリプト言語には固有の役割があり、固有の失敗のしかたがあります。汎用性と誕生当時の事情を踏まえて選ばれた Lua の既定の振る舞いが、その用途にはまさに不向きだったのです。

何度も立ち戻る失敗があります。第4章の中心となる例です。私たちのデモゲーム Lantern Night は、プレイヤーの最高得点をファイルに保存します。Lua 版では、読み込みを次のように書いていました。

\begin{lstlisting}[language=luatwin]
local best = tonumber(lt.load_save("lanternnight_best") or "0")
\end{lstlisting}

\noindent
正しく見えます。レビューも通り、出荷もされました。それでも保存が書き込み途中で中断されると、最初のフレームでクラッシュしました。途中で中断された書き込みは0バイトのファイルを残します。`lt.load_save` は `nil` ではなく空文字列 `""` を返します。Lua では空文字列が\emph{真として扱われる}ため、`or "0"` の代替値は使われません。そして `tonumber("")` は `nil` になり、それ以降の算術式をすべて壊してしまいます。空文字列を真とすること、エラーを `nil` で表すこと、`nil` が黙って伝播すること、静的検査がないこと。この四つの設計判断が重なり、一度のディスク書き込みの中断が、コードを読むだけでは見抜けない最初のフレームのクラッシュに変わりました。

wick では、このバグを書けません。保存専用の例外処理を追加したからではありません。`nil` かもしれない値を使うとき、先に `nil` の場合を扱わなければ、型システムがコンパイルを許さないからです。wick でこの行を書くなら、コンパイラが受け付ける形は次のとおりです。

\begin{lstlisting}[language=wick]
let best = num(lt.load_save("lanternnight_wick_best") ?? "") ?? 0
\end{lstlisting}

\noindent
この一行に、言語の基本思想がすべて詰まっています。以下の設計判断は、どれもこの思想に沿っています。

\section{Lua から学んだこと}

Lua で困った点と、そこから生まれた wick の判断をまとめます。どれも lint ルールやスタイルガイドではありません。コンパイラまたは実行時の仕組みが強制するため、締め切りに追われても忘れようがありません。

\begin{center}
\begin{tabular}{@{}p{0.46\linewidth}p{0.46\linewidth}@{}}
\toprule
\textbf{Lua で困った点} & \textbf{wick の答え} \\
\midrule
実行時の \texttt{nil} エラー &
静的な型と \texttt{T?} のオプショナル値。未確認の使用はコンパイルできない。 \\
タイプミスが暗黙のグローバル変数を作る &
宣言した変数のみ。未宣言の名前への代入はコンパイルエラー。 \\
添字は1から、\texttt{continue} がない &
添字は0からの \texttt{[\,]}、\texttt{continue}、\texttt{len(x)}。 \\
\texttt{""} と \texttt{0} が真になる &
条件は \texttt{bool} に限定。真偽の暗黙変換はない。 \\
フレーム途中の GC 停止 &
回収はフレームの\emph{間}だけで実行。 \\
スタック添字による C バインディング &
エンジン呼び出しの型をコンパイル時に検査。 \\
機種によって \texttt{math.random} が異なる &
\texttt{rand()} は固定の xorshift。どこでも決定的に動作。 \\
iOS では JIT が禁止 &
JIT を使わず、静的な型に応じたバイトコードを生成。 \\
\bottomrule
\end{tabular}
\end{center}

ここには共通するパターンがあります。各行の Lua の振る舞いは、抽象的には合理的でも、ゲームループでは危険です。真偽の暗黙変換は便利ですが、得点が0のときに気づかないまま別の分岐へ進むと困ります。暗黙のグローバル変数は短く書けますが、タイプミスが `nil` になり、三つ先の関数で表面化することがあります。好きなときに動くガベージコレクタはバッチ処理なら問題なくても、アクションゲームではフレーム落ちの原因です。wick の答えはいつも同じです。Lua が実行時に動的に行うことを、コンパイル時に静的に行う。あるいは GC なら、ホストが選んだ時点で行います。

\section{設計の原則}

言語全体を貫く原則は四つあります。後の章で、ある機能がなぜそのように動くのかを説明するとき、答えはほぼ必ずこのいずれかです。

\paragraph{小さいこと。} wick は意図的に Lua 1.0 のような小ささにしています。字句解析器、コンパイラ、仮想マシン、ガベージコレクタ、組み込み関数を含む実装全体が、依存ライブラリのない約二千行の \mbox{C\texttt{++}} です。v0.3 にはクロージャ、クラス、モジュール、メタテーブル、可変長引数、文字列メソッドがありません。省いた機能はすべて、第8章に回避策とロードマップ上の位置づけを記載しています。小ささは、偶然そうなったことへの言い訳ではなく、守ろうとしている特徴です。午後だけで読み切れる言語なら、午前3時にも振る舞いを予測できます。

\paragraph{型があること。} どの式にも静的な型があり、コンパイラがその式のコードを生成する時点で型は分かっています。推論の材料がある限り型を推論するので、書き方は軽いままです。`let x: num = 10` と書かずに `let x = 10` と書けます。それでも検査は全体に及びます。型システムの中心はオプショナル値 `T?` です。「値が存在しないかもしれない」という事実をコンパイラが追跡し、プログラマがその可能性を処理する必要があります。

\paragraph{決定的であること。} 同じプログラムに同じ入力を与えれば、どの機種でも同じフレームを生成します。`rand()` は固定の xorshift 生成器で、時刻によるシード設定はありません。実行ごとに変化をつけたい場合は、自分で選び記録できる数値を `srand()` に渡します。`LANTERN_FIXED_DT` を使うとフレーム時刻の刻みも固定され、ゲームプレイ全体をバイト単位で同じように再現できます。これがスクリーンショットを使った CI を可能にし、第6章ではその上にテストの習慣を築きます。

\paragraph{フレームを意識すること。} wick は毎秒60フレームのループの中で動くことを前提にしています。ガベージコレクタはホストが要求したときだけ動きます。lantern が要求するのは、画像を表示した後、各フレームに一度だけです。停止が画面に現れないタイミングを選んでいます。エンジン API は、ホストが毎フレーム呼ぶ二つの関数 `update(dt)` と `draw()` を中心に構成されています。ゲームにたまたま埋め込まれた汎用ツールではなく、字句解析器の段階からゲームループに合わせた言語です。

\section{本書の読み方}

最初は先頭から順に読み、その後はリファレンスとして使うことを想定しています。

\begin{itemize}
\item \textbf{第2章}はチュートリアルです。エンジンを導入し、最初のプログラムを動かし、小さなランプゲームを段階的に完成させます。厳密なリファレンスの章に進む前に、動くコードを手元に用意できます。
\item \textbf{第3章}は言語リファレンスです。字句構造、型、スコープ、式、優先順位、文、制御フロー、関数を扱います。
\item \textbf{第4章}は wick の核であるオプショナル値を扱います。本章で紹介した0バイトの保存ファイルの実例も交えて説明します。
\item \textbf{第5章}はコレクションであるリストとマップ、そして構造体の代わりに使う並行リストの書き方を扱います。
\item \textbf{第6章}は `lt.*` のエンジンインターフェースを説明します。フレームの契約、型付き API の表、決定性、検証を扱います。
\item \textbf{第7章}は実装を開いて見ます。一回走査の型付きコンパイル、スタック VM、フレーム境界でのガベージコレクション、自分の \mbox{C\texttt{++}} への wick の組み込み方を説明します。
\item \textbf{第8章}は設計の根拠です。同じ実際のゲームで wick と Lua を比べ、率直なトレードオフとロードマップの考え方を述べます。
\item \textbf{付録}には EBNF による完全な文法、コンパイルエラーの全一覧、組み込み関数とエンジン API の早見表があります。
\end{itemize}

本書のコードは、別の言語だと明記されない限り、\texttt{この書体}で示した wick のコードです。どのコード例も実際の v0.3 の wick で、コンパイルできるものです。コンパイルエラーを示すため意図的に間違えた例には、対応するエラーメッセージを添えています。Lua の例は wick の例に対応するもので、Lua と明記しています。

率直さについて一つ。本書では、今の wick より Lua の方が優れている点を何度か述べます。動的なテーブルは入れ子のデータをテーブルとして保持できます。複数の戻り値は実際に不足しています。Lua のエコシステムは巨大で、こちらは一つのエンジンです。対象を褒めるだけの言語の本は、あまり役に立ちません。wick が存在する価値は、特定の、しかも高くつく種類のバグを取り除くことにあります。何を取り除き、何を取り除かないのかは、正確に述べるべきです。
